\documentclass[11pt]{article}
\usepackage[margin=1in]{geometry}
\usepackage{amsmath,amssymb,amsthm,booktabs,hyperref}
\newtheorem{theorem}{Theorem}
\newtheorem{lemma}{Lemma}
\newcommand{\ospt}{\operatorname{ospt}}
\title{An effective proof of the strongly unimodal rank-difference conjecture}
\author{Principia Math in Progress}
\date{12 September 2026}
\begin{document}
\maketitle

\paragraph{Task 1. Resolve the rank-difference conjecture.}
A strongly unimodal sequence increases strictly to a single peak and then
decreases strictly. Its rank is the number of entries after the peak minus
the number before it. Let \(u(m,n)\) count such positive integer sequences
of total weight \(n\) and rank \(m\), let \(p(n)\) count partitions of \(n\),
and let \(N(0,n)\) count partitions with largest part equal to their length.
Zang's Conjecture 5.1 \cite{Zang}, recorded as MathDB 364074, asks for
\[
 u(0,n)-u(1,n)\geq \tfrac12N(0,n)\qquad(n\geq8).
\]

\begin{theorem}
The conjectured inequality holds for every \(n\geq8\), with equality
if and only if \(n=11\).
\end{theorem}

\paragraph{Reduction and exact kernel.}
Zang's identities (4.6)--(4.7) give
\[
 u(0,n)+u(1,n)=\tfrac12(p(n)+N(0,n))\qquad(n\geq1).
\]
Thus the required deficit, multiplied by two, is
\(E_n=p(n)-4u(1,n)\). Set \(P(q)=(q;q)_\infty^{-1}\),
\(T_r=r(r+1)/2\), and \(U_1(q)=\sum_{n\geq1}u(1,n)q^n\).
The fixed-rank identity \cite[Theorem 1.1]{Peak}, with rank one, gives
\begin{equation}\label{U}
 (q;q)_\infty U_1(q)=
 \sum_{r\geq2}(-1)^r\sum_{j=0}^{r-2}q^{T_r+jr}.
\end{equation}
Indeed, its summation variable \(k\geq1\) has prefactor
\(q^{1+k(k+1)/2+k}=q^{T_{k+1}}\); expanding
\((q^{k(k+1)}-1)/(1-q^{k+1})\) and putting \(r=k+1\) gives
\eqref{U}.
Writing \(E(q)=\sum_{n\geq0}E_nq^n=P(q)-4U_1(q)\), we obtain
\[
 E(q)=P(q)K(q),\qquad
 K(q)=1+4\sum_{r\geq2}(-1)^{r+1}\sum_{j=0}^{r-2}q^{T_r+jr}.
\]
The two terms \(j=r-1,r\) have the two pentagonal exponents
\(r(3r-1)/2,r(3r+1)/2\). Euler's pentagonal identity therefore yields
\[
 K(q)=4(q;q)_\infty+H(q),\qquad
 H(q)=-3+4\sum_{r\geq1}(-1)^{r+1}\sum_{j=0}^r q^{T_r+jr}.
\]
Consequently \(E_n=[q^n]P(q)H(q)\) for \(n\geq1\).
For \(\Re z>0\), putting \(r=a+b,\ j=a\) gives
\begin{equation}\label{theta}
 H(e^{-z})=1-4\sum_{a,b\geq0}(-1)^{a+b}e^{-zQ(a,b)},\quad
 Q(a,b)=\tfrac32a^2+2ab+\tfrac12b^2+\tfrac12(a+b).
\end{equation}
All these series converge absolutely in the asserted domains.

\paragraph{A uniform local estimate.}
Set
\[
 x_0=\frac{2\pi}{\sqrt{239999}},\qquad
 \rho=\sqrt2\,x_0,\qquad h_0=\sqrt\rho .
\]
We will establish the effective bound
\begin{equation}\label{local}
 \left|H(e^{-z})-\tfrac3{16}z^2\right|
 <0.041\,\tfrac3{16}|z|^2
 \quad(0<|z|\leq\rho,\ |\arg z|\leq\pi/4).
\end{equation}
This is a computer-assisted estimate with an explicit finite certificate,
described next; it is not an asymptotic with an unspecified error.

Write \(z=h^2\omega\), with \(h>0,\ |\omega|=1\), and set
\[
 f(X,Y)=\exp[-\omega(3X^2/2+2XY+Y^2/2+h(X+Y)/2)].
\]
Apply order-16 Boole summation \cite[24.17.1--2]{DLMF} in each variable.
In its half-line scaled form it is
\begin{equation}\label{boole}
 \sum_{m\geq0}(-1)^m g(mh)
 =\sum_{d\in D}c_dh^dg^{(d)}(0)+R(g),\qquad
 |R(g)|\leq C h^{15}\int_0^\infty|g^{(16)}(t)|\,dt,
\end{equation}
where
\[
 D=\{0,1,3,5,\ldots,15\},\quad c_0=\tfrac12,\quad
 c_{2k-1}=-\frac{(4^k-1)B_{2k}}{(2k)!},\quad C=\frac4{\pi^{16}}.
\]
For completeness, the exact remainder kernel has norm at most
\(2\pi^{-16}\sum_{k\geq0}(2k+1)^{-16}<4\pi^{-16}\), by the
Euler-polynomial Fourier series \cite[24.8.5]{DLMF}. Scaling its integral
explains \(h^{15}\). Gaussian decay permits passage to the half-line.
The same formula applies to real and imaginary parts, or directly to a
complex integrand using the real remainder kernel.

Here are full specifications for reproducing the finite bound.
Let \(P_{ab}=f^{-1}\partial_X^a\partial_Y^bf\). These rational polynomials
in \(X,Y,h,\omega\) are generated by
\[
 P_{00}=1,\quad
 P_{a+1,b}=\partial_XP_{ab}-\omega(3X+2Y+h/2)P_{ab},\quad
 P_{a,b+1}=\partial_YP_{ab}-\omega(2X+Y+h/2)P_{ab}.
\]
On the nonnegative quadrant,
\[
 |f(X,Y)|\leq e^{-\alpha X^2-\beta Y^2},\quad
 \alpha=\frac3{2\sqrt2},\quad\beta=\frac1{2\sqrt2}.
\]
For a polynomial \(P=\sum p_{ijk\ell}X^iY^jh^k\omega^\ell\), define
\begin{align*}
 I_{y=0}(P)&=\sum_{i,k,\ell}|p_{i0k\ell}|h_0^kG_i(\alpha),\\
 I_{x=0}(P)&=\sum_{j,k,\ell}|p_{0jk\ell}|h_0^kG_j(\beta),\\
 I_2(P)&=\sum_{i,j,k,\ell}|p_{ijk\ell}|h_0^kG_i(\alpha)G_j(\beta),\\
 G_i(c)&=\frac{\Gamma((i+1)/2)}{2c^{(i+1)/2}}.
\end{align*}
They bound the corresponding absolute derivative integrals, uniformly
over \(h\leq h_0\) and the sector in \eqref{local}.
The three remainder terms \(R_XL_Y,L_XR_Y,R_XR_Y\) are bounded by
\[
 R_X^*=Ch_0^{15}\sum_{d\in D}|c_d|h_0^dI_{y=0}(P_{16,d}),\quad
 R_Y^*=Ch_0^{15}\sum_{d\in D}|c_d|h_0^dI_{x=0}(P_{d,16}),\quad
 R_{XY}^*=C^2h_0^{30}I_2(P_{16,16}).
\]
Fubini is justified by the same integrable Gaussian derivative bounds.
The finite boundary term for \(H\) is
\[
 B_H(z)=1-4\sum_{a,b\in D}c_ac_bh^{a+b}P_{ab}(0,0).
\]
Every surviving monomial satisfies \(a+b+k=2\ell\), so this is a
polynomial in \(z=h^2\omega\). Its constant and linear terms vanish, and
its next coefficients are exactly
\[
 B_H(z)=\tfrac3{16}z^2+\tfrac38z^3+\tfrac{227}{256}z^4
             +\tfrac{41}{16}z^5+\tfrac{181499}{20480}z^6+\cdots .
\]
The attached verifier uses exact rational arithmetic for all polynomials
and 220-bit Arb balls for the Gaussian moments and endpoints. It certifies:
\begin{center}
\begin{tabular}{lr}
\toprule
Quantity & Certified bound\\
\midrule
\(R_X^*\) & \(<4.357\,10^{-8}\)\\
\(R_Y^*\) & \(<1.760\,10^{-11}\)\\
\(R_{XY}^*\) & \(<9.583\,10^{-10}\)\\
\(4(R_X^*+R_Y^*+R_{XY}^*)\) & \(<1.79\,10^{-7}\)\\
\(\sum_{j\geq3}|[z^j]B_H|\,\rho^j\) & \(<2.34\,10^{-6}\)\\
\((3/16)\rho^2\) & \(>6.168\,10^{-5}\)\\
Total error divided by the leading term & \(<0.040808\)\\
\bottomrule
\end{tabular}
\end{center}
The actual enclosed ratio is \(0.040807231163328\ldots\).
After division by \(h^4\), each polynomial-tail term has degree at least
two in \(h\), each single remainder bound degree at least eleven, and
each double remainder bound degree at least twenty-six. All their
coefficients are nonnegative. Their endpoint bounds therefore hold for
every \(0<h\leq h_0\), proving \eqref{local}.

\paragraph{The major arc.}
For \(n\geq10000\), put \(A=\pi^2/6,\ N=n-1/24,\ x=\sqrt{A/N}\leq x_0\).
Cauchy's coefficient formula is
\[
 E_n=\frac1{2\pi}\int_{-\pi}^{\pi}
 P(e^{-x-iy})H(e^{-x-iy})e^{n(x+iy)}\,dy.
\]
The usual eta transformation, with the square root in the right half-plane,
is the exact identity
\begin{equation}\label{eta}
 P(e^{-z})=\sqrt{\frac z{2\pi}}e^{A/z-z/24}
          P(e^{-4\pi^2/z}).
\end{equation}
On \(|y|\leq x\), \(\Re(1/z)\geq1/(2x)\).
If \(|w|\leq r=e^{-2\pi^2/x}\), the product series gives
\(|P(w)-1|\leq \exp(r/(1-r)^2)-1<3r\);
here \(2\pi^2/x_0>1539\).
Multiplying this error with the actual enclosed local error above still
gives a combined relative error below \(0.041\).

Set \(y=x^{3/2}s\). Factoring the leading contribution yields
\[
 \text{major arc}=\frac{(3/16)x^4e^{2A/x}}{2\pi\sqrt{2\pi}}J,
\]
where the unperturbed amplitude and phase are
\[
 W_x(s)=(1+xs^2)^{5/4}e^{-As^2/(1+xs^2)},\qquad
 \Phi_x(s)=\frac{A\sqrt{x}s^3}{1+xs^2}
                     +\frac52\arctan(\sqrt{x}s),
\]
and \(|s|\leq1/\sqrt{x}\). The relative error in this integrand is less
than \(0.041\), possibly complex.
On \(0\leq s\leq1\), \(W_x(s)\geq e^{-A}\) and
\(\Phi_x(s)\leq(A+5/2)\sqrt{x_0}<0.47\). Thus the real half-integral
on this interval exceeds \(e^{-A}\cos(0.47)>0.172\).
For \(1\leq s\leq3/2\),
\[
 \Phi_x(s)\leq ((27/8)A+15/4)\sqrt{x_0}<\pi/2,
\]
so that contribution is nonnegative.

For \(3/2\leq s\leq1/(2\sqrt{x})\),
\(W_x(s)\leq(5/4)^{5/4}e^{-4As^2/5}\).
For \(1/(2\sqrt{x})\leq s\leq1/\sqrt{x}\),
\(W_x(s)\leq2^{5/4}e^{-A/(5x)}\).
Writing \(c=4A/5\), the remaining half-tail is consequently bounded by
\[
 (5/4)^{5/4}\frac{e^{-9c/4}}{3c}
 +\frac{2^{5/4}}{2\sqrt{x}}e^{-A/(5x)}<0.0175.
\]
The full absolute half-integral is bounded by
\[
 (5/4)^{5/4}\frac{\sqrt\pi}{2\sqrt c}
 +\frac{2^{5/4}}{2\sqrt{x}}e^{-A/(5x)}<1.022.
\]
The \(x\)-dependent terms increase on \(0<x\leq x_0\);
the remaining endpoint inequalities are certified by Arb.
Conjugate symmetry and the relative error bound now give
\[
 J>2(0.172-0.0175-0.041\cdot1.022)>0.22 .
\]
Thus the major arc exceeds the positive quantity
\begin{equation}\label{M}
 M_n=\frac{(3/16)\,0.22}{2\pi\sqrt{2\pi}}x^4e^{2A/x}.
\end{equation}

\paragraph{The minor arc.}
For \(a=e^{-x}\) and \(x\leq|y|\leq\pi\), keeping just the first
logarithmic layer of the Euler product gives
\begin{align*}
 \log P(e^{-x})-\log|P(e^{-x-iy})|
 &\geq\sum_{m\geq1}a^m(1-\cos(my))\\
 &=\frac{a(1+a)(1-\cos y)}
 {(1-a)((1-a)^2+2a(1-\cos y))}.
\end{align*}
The last expression increases with \(1-\cos y\).
At \(|y|=x\), use \(a\geq1-x\), \(1-a\leq x\), and
\(x^2/2-x^4/24\leq1-\cos x\leq x^2/2\). The expression is at least
\[
 \frac{(1-x_0)(2-x_0)(1/2-x_0^2/24)}{2x}>\frac{0.49}{x}.
\]
Also the series defining \(H\) gives, since \(T_r\geq r^2/2\),
\[
 |H(e^{-x-iy})|
 \leq3+4\sum_{r\geq1}(r+1)e^{-xr^2/2}
 \leq3+\frac4x+8\sqrt{\frac\pi{2x}}<\frac{5.174}{x}.
\]
For the sum estimate, compare each summand on \([r-1,r]\) with
\((t+2)e^{-xt^2/2}\) and integrate.
Using the real form of \eqref{eta}, the minor-arc length at most \(2\pi\),
and \eqref{M}, we obtain
\[
 \frac{|\text{minor arc}|}{M_n}
 \leq (1+3e^{-4\pi^2/x})
 \frac{2\pi\,5.174}{(3/16)\,0.22}\,
 x^{-9/2}e^{-0.49/x}<7\,10^{-6}.
\]
Since \(x_0<0.49/(9/2)\), both factors increase throughout the interval.
Their endpoint product is below \(0.000006577\), including the eta residual.
Thus \(E_n>0\) for every \(n\geq10000\).

\paragraph{The remaining integers and equality.}
For \(k_n=[q^n]K(q)\), Euler's identity gives the exact recurrence
\[
 E_n=k_n-\sum_{r\geq1}(-1)^r
 \left(E_{n-r(3r-1)/2}+E_{n-r(3r+1)/2}\right),
 \quad E_j=0\ (j<0).
\]
The attached integer certificate evaluates it for \(0\leq n\leq10000\).
It verifies \(E_n\geq0\) for all 9,993 integers \(8\leq n\leq10000\),
and its only zero there is \(11\). This finite interval overlaps the
analytic interval, proving the theorem and its equality assertion.
Independent direct enumeration checks the generating-function reduction
through \(n=120\), and a separate enumeration of 215,307 ordinary
partitions checks both rank identities through \(n=40\).

\medskip\noindent
The following tasks are proposed after the full proof of Task 1.

\paragraph{Task 2. Find a combinatorial proof.}
Construct an injection from four labelled copies of the rank-one strongly
unimodal sequences of weight \(n\) into the partitions of \(n\), for
\(n\geq8\), and describe the complement, including its disappearance at \(11\).

\paragraph{Task 3. Prove a positive generating-function formula.}
Express \(\sum_{n\geq8}E_nq^n\) in a manifestly nonnegative form,
with a combinatorial interpretation for its coefficients.

\paragraph{Task 4. Quantify the positive deficit.}
Derive simple explicit two-sided bounds for \(E_n/p(n)\), identify
the leading constant in its order-\(n^{-1}\) decay, and determine sharp
finite thresholds for those bounds.

\paragraph{Task 5. Extend the effective method to varying ranks.}
Obtain effective estimates uniform in both \(m\) and \(n\) for
\(u(m,n)-u(m+1,n)\), and investigate which range of \(m\) permits a
full log-concavity proof for the rank distribution.

\paragraph{Verification and attribution.}
The proof uses Zang's rank identity, the earlier fixed-rank generating
function, and classical Boole and eta transformations. These are prior
results, not new claims. The new certificate supplies explicit positivity
and the full finite range for the target conjecture.
Run \texttt{python} \path{problems/364074/verify.py} from the campaign directory
with python-flint 0.9.0. Its exact integers/rationals and 220-bit Arb
enclosures are part of this computer-assisted proof. The wrapper rejects
optimized Python, which would disable the historical certificate's assertions.
The recorded run passed all checks in about 17 seconds. No proof-assistant
formalization is claimed.

\begin{thebibliography}{9}
\raggedright
\bibitem{Zang}
W.~J.~T.~Zang, \emph{Unimodality of the Rank on Strongly Unimodal Sequences},
\href{https://arxiv.org/abs/2407.18186}{arXiv:2407.18186v1} (2024),
Conjecture 5.1 and equations (4.6)--(4.7).
\bibitem{Peak}
K.~Bringmann, C.~Jennings-Shaffer, K.~Mahlburg and R.~Rhoades,
\emph{Peak positions of strongly unimodal sequences},
Trans. Amer. Math. Soc. \textbf{371} (2019), 7087--7109,
\href{https://doi.org/10.1090/tran/7791}{doi:10.1090/tran/7791};
\href{https://arxiv.org/abs/1806.03217}{arXiv:1806.03217v2}.
\bibitem{DLMF}
NIST Digital Library of Mathematical Functions,
\href{https://dlmf.nist.gov/24.17}{24.17.1--2} and
\href{https://dlmf.nist.gov/24.8.E5}{24.8.5}.
\end{thebibliography}
\end{document}
